
Un modelo neuronal es una formalización matemática del comportamiento de una neurona real, que trata de imitarla, estudiarla o comunicarse con ella. En su acepción más habitual es la modelización de la generación de potenciales de acción o \textit{spikes} \cite{Fortuna2023}. Este tipo de señales, generan una serie de picos, que hacen que las poblaciones neuronales se muevan entre estados. Dichos estados son lo que le da sentido a la comunicación neuronal \cite{Rabinovich2008}.

Algunos de estos modelos, presentan un sistema dinámico con un régimen caótico, dicho régimen presenta oscilaciones, en los que hay un estado estable y uno inestable dentro de su espacio fásico \cite{Rulkov2002}. Debido a esto, para poder representarlos, se usan ecuaciones diferenciales.  
